% !TeX root = ../../Book2.tex
\OptionalSection{现代 Fourier 限制问题}{b2:sec:2306}

本节区分已完整证明的 Stein–Tomas 范围、必要条件给出的广泛目标和现代几何方法。研究进展按正式出版的具体结果引用，不把一个方法名称当作整个猜想的证明。

\subsection{限制猜想}
\label{b2:subsec:230601}

对紧非零曲率超曲面，一种经典延拓目标是
\[
 E:L^s(S)\to L^q(\mathbb R^d),\qquad
 q>\frac{2d}{d-1},\quad
 \frac{d+1}{q}\le\frac{d-1}{s'}.
\]
必要条件已在本章证明；称为限制猜想时，这些不等式描述的是目标范围，不是本章已经取得的定理。具体维数、曲面及端点版本需要分别陈述。本书采用这一经典目标作为组织方法的背景，不以未核验的“最新范围”替代准确的定理。

\ProseParagraphBreak
当 \(s=2\)，Stein–Tomas 已达到帽必要条件边界。进一步的困难主要在输入指数改变后的范围，以及许多方向相互作用的控制。球面、抛物面与锥面有不同对称性和曲率结构，不能无说明地共用同一端点命题。

\subsection{双线性归约}
\label{b2:subsec:230602}

把 \(g\) 分成小帽后，\(|Eg|^2\) 是不同帽延拓的乘积之和。
相邻帽的乘积可在更细尺度继续分解；分离帽的法向有角度，则有额外双线性增益。这形成双线性归约的基本思路。
归约的完整形式需控制近对角项、远对角项、尺度求和及输出空间的准范数，不是简单地丢弃相邻帽。

\ProseParagraphBreak
\chapref{b2:ch:25}将完整计算二维抛物线上分离区间的乘积 \(L^2\) 估计：和映射的 Jacobi 行列式恰好等于法向分离量。这个计算是现代双线性思想的一个可独立核验模型。一般高维双线性限制理论需要更复杂的几何，不作为本章核心证明依赖。

\subsection{波包}
\label{b2:subsec:230603}

一个半径 \(R^{-1/2}\) 的曲面帽产生切向宽度 \(R^{1/2}\)、法向长度 \(R\) 的波包。先在帽内做频率局部化，再按切向空间位置分解输入，可把延拓近似组织成许多这样的管。
频率宽度与空间宽度互为倒数，法向长度则来自曲率的二次变化，这与 Knapp 管的计算一致。

\ProseParagraphBreak
波包不是严格紧支撑于一条管：光滑截断提供的是快速衰减尾部。将波包当作管特征函数会丢失相位、正交性和尾部误差。完整方法必须同时记录系数平方和、管的方向及误差求和。

\subsection{解耦}
\label{b2:subsec:230604}

解耦比较总延拓的 \(L^q\) 范数与各个小帽延拓范数的平方和。
它允许各帽内部保持相干，再控制帽与帽之间的组合；因此不同于只对输入使用 Plancherel。
抛物面的一个正式出版结果见 \cite{BourgainDemeter2015Proof}，系统处理见 \cite{Demeter2020Fourier}；\chapref{b2:ch:25}将准确写出其带尺度损失的模型版本并标注深层输入。

\ProseParagraphBreak
不同曲面、不同切分尺度和不同 \(q\) 对应不同解耦常数。必须保留加权空间球、帽大小与 \(R^\epsilon\) 因子，才能在尺度归纳中使用。把解耦粗略概括成“帽之间正交”会把 \(L^2\) 的等距事实与非 Hilbert 空间的深层估计混淆。

\subsection{多项式分割}
\label{b2:subsec:230605}

多项式分割把空间中的质量分配到许多胞腔，并把靠近多项式零集的波包单独处理。横穿胞腔的管可利用交点次数，贴近零集的管则需低维几何或横截估计。这两部分的不同控制是方法的核心。
相关正式出版的限制应用见 \cite{Guth2016Restriction}。本章只说明方法处理什么障碍，不在没有完整胞腔计数、切向管分析和尺度递推的情况下宣称一个新范围。

\ProseParagraphBreak
现代限制理论的共同组织原则是先辨认局部相干的波包，再利用方向之间的几何分离。驻相、平方函数、横截性和尺度归纳分别处理不同环节。阅读具体定理时应先核对曲面、输入指数、输出指数、局部半径和允许的尺度损失，再判断它能否用于目标问题。
